\section{Discussion}
\label{sec:discussion}

\paragraph{Why pair-disjoint calibration matters.}
The TempGlitch follow-up improves on the earlier identical-episode family by preventing the
calibration normals from sharing source, suspected pair, or episode identity with evaluation
episodes. This matters because paired normal and buggy captures may share scene layout, assets,
camera motion, or recording conditions. A threshold calibrated on a related normal could encode
pair-specific information rather than a transferable notion of normality. The repaired split does
not eliminate every possible dataset bias, but it closes the documented cross-role overlap.

\paragraph{What the main result says.}
The best recorded LeWM row has higher AUROC and AUPRC than the strongest simple baseline row on
the exact same frozen TempGlitch support. Its thresholded F1 is also substantially higher. At the
same time, overlapping AUROC intervals, two-episode calibration, 12 normal negatives, and
FPR@95TPR of 0.7500 make the observation fragile. The scientifically useful conclusion is that
latent surprise warrants further controlled study, not that it has won a general benchmark.

\paragraph{Why R5-WOB is not a binary benchmark.}
The R5-WOB bundle validates execution and records positive-class behavior under a threshold
calibrated from normal episodes. Because its evaluation role has no normal-negative episodes,
AUROC and false-positive behavior cannot be estimated as a binary benchmark. R5-XGame adds held
normal negatives and therefore supports bounded binary metrics, but it remains a separate,
positive-heavy non-locked family.

\paragraph{Future comparisons.}
The current simple baselines are useful controls but do not exhaust plausible alternatives.
Stronger evidence would require an exact-support learned video baseline, additional seeds,
predeclared SIGReg and action-conditioning comparisons, and broader datasets with verified
temporal annotations. Those experiments require a new compute phase and are not inferred from the
existing artifacts.

\begin{table}[t]
\caption{Evidence surfaces and their claim boundaries.}
\label{tab:limitations-claim-boundary}
\centering
\small
\begin{tabularx}{\textwidth}{@{}lXX@{}}
\toprule
Surface & Supports & Does not support \\
\midrule
TempGlitch follow-up & Bounded same-support separation on one frozen pair-disjoint non-locked split & Broad superiority, temporal localization, state of the art \\
R5-XGame & Bounded ranking on one positive-heavy non-locked split & Cross-game generalization or action-conditioning benefit \\
GlitchBench K2 & Bounded image-level ranking on one synthetic-normal split & Temporal localization, natural-normal gameplay, or broad superiority \\
R5-WOB & Pipeline execution and positive-probe signal presence & Binary discrimination or false-positive performance \\
Locked test & No result; gate remains closed & Any locked-test metric or conclusion \\
\bottomrule
\end{tabularx}
\end{table}

\begin{table}[t]
\caption{Reviewer risks and the evidence-based response used in this draft.}
\label{tab:reviewer-risk}
\centering
\small
\begin{tabularx}{\textwidth}{@{}lXX@{}}
\toprule
Risk & Evidence symptom & Paper treatment \\
\midrule
Small negative support & 12 normal negatives in each binary family & Report counts and confidence intervals beside metrics \\
Calibration fragility & TempGlitch threshold uses two normal episodes & Name calibration episodes and avoid deployment claims \\
High false-positive rate & TempGlitch FPR@95TPR is 0.7500 & Report explicitly rather than relying on AUROC alone \\
Seed sensitivity & R5-XGame seed44 is strongest & Treat best-row ranking as split-bounded \\
Binary label granularity & No verified temporal spans & Forbid temporal-localization wording \\
Protocol mismatch & Zero-action TempGlitch and real-action XGame & Keep result families separate \\
Public non-locked evidence & No locked-test result & Avoid final-benchmark language \\
\bottomrule
\end{tabularx}
\end{table}

